\chapter{Tổng kết và hướng phát triển}
\label{ch:tong-ket-va-huong-phat-trien}

\section{Kết quả đạt được}
\label{sec:tk-ket-qua}

Nhóm đã xây dựng hệ thống CoSpace gồm ứng dụng web React, máy chủ Spring Boot và cơ sở dữ liệu PostgreSQL trên Supabase, phục vụ bốn nhóm người dùng: khách hàng, nhân viên, quản lý chi nhánh và quản trị viên hệ thống. Bảng~\ref{tab:tk-so-lieu} tóm tắt quy mô sản phẩm; các số liệu được đếm trực tiếp từ mã nguồn.

\begin{table}[H]
\centering
\caption{Quy mô sản phẩm}
\label{tab:tk-so-lieu}
\small
\renewcommand{\arraystretch}{1.2}
\begin{tabularx}{\linewidth}{|L|P{3.2cm}|}
\hline
\textbf{Hạng mục} & \textbf{Số lượng} \\
\hline
Yêu cầu chức năng được đặc tả (Chương~\ref{ch:yeucau}) & 34 \\
\hline
Bảng dữ liệu (sau migration Flyway \mt{V1}--\mt{V3}) & 32 \\
\hline
Controller REST phía máy chủ & 28 \\
\hline
Điểm cuối API & 153 \\
\hline
Trang chức năng trên giao diện (theo bốn vai trò và trang công khai) & 31 \\
\hline
Ca kiểm thử tự động phía máy chủ & 307 (306 đạt, 1 bỏ qua) \\
\hline
Ca kiểm thử ràng buộc trên PostgreSQL & 17 (17 đạt) \\
\hline
\end{tabularx}
\end{table}

Bảng~\ref{tab:tk-muc-tieu} đối chiếu kết quả với các mục tiêu đã đặt ra ở Mục~\ref{sec:muctieu}.

\begin{table}[H]
\centering
\caption{Đối chiếu kết quả với mục tiêu}
\label{tab:tk-muc-tieu}
\footnotesize
\renewcommand{\arraystretch}{1.2}
\begin{tabularx}{\linewidth}{|P{3.6cm}|L|P{1.9cm}|}
\hline
\textbf{Mục tiêu} & \textbf{Kết quả} & \textbf{Mức độ} \\
\hline
\multicolumn{3}{|l|}{\textit{Mục tiêu nghiệp vụ}} \\
\hline
Số hóa quy trình đặt chỗ & Tìm chỗ trên sơ đồ tầng, giữ chỗ 15 phút, thanh toán VietQR (PayOS) hoặc MoMo, nhận mã đặt chỗ và mã QR để check-in & Đạt \\
\hline
Quản lý nhiều chi nhánh & Phân cấp chi nhánh -- tầng -- không gian; bảng giá, dịch vụ, chính sách hủy riêng cho chi nhánh; dữ liệu mẫu 3 chi nhánh, 9 tầng & Đạt \\
\hline
Tự động hóa hủy và hoàn tiền & Tính tiền hoàn theo bậc chính sách, giải phóng chỗ ngay khi hủy; việc trả tiền cho khách vẫn làm thủ công rồi đánh dấu trên hệ thống, hoặc hoàn bằng voucher & Đạt một phần \\
\hline
Hỗ trợ vận hành quầy & Đơn tại quầy, thu tiền mặt, check-in, check-out, dịch vụ phát sinh, gia hạn, phí trả chỗ muộn & Đạt \\
\hline
Kết nối cộng đồng & Hồ sơ năng lực, gợi ý đối tác theo độ tương đồng Jaccard, lời mời kết nối, bảng tin có thẻ chủ đề & Đạt \\
\hline
Báo cáo quản trị & Doanh thu, số đơn, tỷ lệ hủy, so sánh chi nhánh, xuất tệp & Đạt \\
\hline
\multicolumn{3}{|l|}{\textit{Mục tiêu kỹ thuật}} \\
\hline
Kiến trúc phân tầng & Controller -- service -- repository phía máy chủ; giao diện tổ chức theo vai trò & Đạt \\
\hline
Chống đặt trùng lịch & Khóa tư vấn \mt{pg\_advisory\_xact\_lock} và ràng buộc \mt{EXCLUDE USING gist}; đã kiểm thử ở mức đơn vị và trên PostgreSQL, chưa kiểm thử tải đồng thời qua API & Đạt \\
\hline
Xử lý webhook an toàn & Kiểm tra chữ ký, bỏ qua webhook lặp, xử lý thanh toán muộn và thanh toán trùng. Khóa MoMo đang là khóa sandbox & Đạt một phần \\
\hline
Bảo mật và phân quyền & Hai loại JWT, phân quyền theo vai trò, \mt{BranchAccessGuard}; còn các điểm tồn đọng ở Mục~\ref{sec:tk-han-che} & Đạt một phần \\
\hline
Trợ lý ảo có kiểm soát & 9 công cụ gọi hàm, thao tác ghi dữ liệu phải qua bước xác nhận & Đạt \\
\hline
Triển khai đám mây & Máy chủ đóng gói Docker chạy trên Render, giao diện trên Vercel, cơ sở dữ liệu trên Supabase & Đạt \\
\hline
\end{tabularx}
\end{table}

Về kỹ thuật, những phần nhóm đã làm và kiểm chứng được gồm:
\begin{itemize}
    \item \textbf{Chống đặt trùng hai lớp:} khóa tư vấn làm các giao dịch đặt cùng một chỗ chạy lần lượt; ràng buộc loại trừ trong cơ sở dữ liệu là lớp chặn cuối cùng. Luồng tạo lịch bảo trì dùng chung khóa với luồng đặt chỗ.
    \item \textbf{Tính toán phía máy chủ:} giá, số đơn vị thời gian, chi nhánh của đơn và tiền dịch vụ đều do máy chủ tính; giao diện chỉ hiển thị báo giá lấy từ API.
    \item \textbf{Máy trạng thái đơn đặt chỗ:} mọi chuyển trạng thái đi qua \mt{BookingStateMachine}; các luồng hủy, check-in và tác vụ định kỳ đều khóa bản ghi và kiểm tra lại trạng thái trước khi ghi.
    \item \textbf{Xử lý thanh toán lặp và muộn:} webhook đến lần hai không được xử lý lại; tiền về sau khi đơn hết hạn hoặc bị hủy được đưa vào hàng chờ hoàn tiền thay vì làm sống lại đơn.
    \item \textbf{Phân quyền đọc lại từ cơ sở dữ liệu:} vai trò và trạng thái tài khoản được nạp lại ở mỗi yêu cầu, nên tài khoản bị khóa mất quyền truy cập ngay.
    \item \textbf{Lược đồ quản lý bằng Flyway:} một cơ sở dữ liệu trống dựng được bằng ba migration \mt{V1}--\mt{V3}.
\end{itemize}

\section{Hạn chế}
\label{sec:tk-han-che}

Hệ thống còn các hạn chế sau. Một số mục đã được ghi trong tài liệu rà soát logic của nhóm và chưa được khắc phục.

\subsection{Bảo mật}
\begin{enumerate}
    \item Khóa bí mật MoMo trong \mt{application.yml} là khóa sandbox công khai và chưa có cơ chế chặn khi chạy ở môi trường production (PayOS đã có \mt{PayosProductionGuard}, MoMo chưa có). Với cấu hình này, người biết khóa có thể giả mạo thông báo thanh toán MoMo.
    \item Token truy cập và token làm mới được lưu trong \mt{localStorage}; chưa có cơ chế thu hồi token, nên khi đăng xuất token cũ vẫn dùng được đến khi hết hạn (tài khoản bị khóa thì bị chặn ngay vì trạng thái được đọc lại từ cơ sở dữ liệu).
    \item API tạo chính sách hủy nhận trực tiếp đối tượng thực thể từ yêu cầu (mass assignment), nên quản lý chi nhánh có thể gửi \mt{id} của một chính sách khác và ghi đè nó.
    \item Nhân viên và quản lý chi nhánh xem được danh sách người dùng của mọi chi nhánh.
    \item Chưa có giới hạn tần suất yêu cầu (rate limiting) cho các API đăng nhập và trợ lý ảo.
    \item Quy tắc cho phép truy cập \mt{/h2-console/**} vẫn còn trong cấu hình bảo mật dù dự án không dùng H2.
\end{enumerate}

\subsection{Nghiệp vụ}
\begin{enumerate}
    \item Chính sách hủy chưa lọc theo loại không gian (\mt{workspace\_type\_id}) và thời gian hiệu lực (\mt{effective\_from}, \mt{effective\_to}) dù các cột này có trong bảng. Các dòng chính sách kiểu cũ (\mt{HOURS\_BEFORE\_START}) vẫn hiển thị ở màn hình quản trị nhưng không được dùng khi tính tiền hoàn.
    \item Hoàn tiền chưa tích hợp với cổng thanh toán: quản lý chi nhánh chuyển tiền cho khách ngoài hệ thống rồi đánh dấu đã xử lý.
    \item Màn hình đặt chỗ tại quầy có hiển thị các lựa chọn MoMo và thẻ, nhưng hệ thống chỉ ghi nhận thanh toán tiền mặt cho đơn tại quầy.
    \item Nhật ký kiểm toán chưa ghi các thao tác đổi vai trò người dùng, xác nhận thu tiền mặt và check-in. Bảng \mt{payment\_events} có trong lược đồ nhưng chưa được sử dụng.
\end{enumerate}

\subsection{Kỹ thuật và vận hành}
\begin{enumerate}
    \item Hai cột \mt{svg\_content} và \mt{layout\_json} của bảng \mt{floors} có trong cơ sở dữ liệu đang chạy nhưng chưa có trong migration \mt{V1\_\_baseline.sql}; một cơ sở dữ liệu mới dựng hoàn toàn bằng Flyway sẽ thiếu hai cột này và chức năng sơ đồ tầng sẽ lỗi.
    \item Các tác vụ định kỳ (hết hạn đơn mỗi 30 giây, cập nhật vòng đời đơn mỗi 5 phút) chỉ chạy đúng khi có một phiên bản máy chủ; chưa có khóa phân tán kiểu ShedLock để chạy nhiều phiên bản.
    \item Bộ nhớ đệm Caffeine nằm trong bộ nhớ của từng phiên bản máy chủ, không chia sẻ giữa các phiên bản.
    \item Máy chủ chạy ở gói miễn phí của Render nên bị tạm dừng khi không có truy cập; yêu cầu đầu tiên sau đó phải chờ máy chủ khởi động lại.
    \item Chưa có quy trình CI tự động chạy kiểm thử; chưa có kiểm thử tích hợp với PostgreSQL thật, kiểm thử tải đồng thời qua API và kiểm thử tự động cho giao diện (Mục~\ref{sec:kt-danh-gia-chung}).
\end{enumerate}

\section{Hướng phát triển}
\label{sec:tk-huong-phat-trien}

\subsection{Ngắn hạn}
Các việc sau khắc phục trực tiếp những hạn chế ở Mục~\ref{sec:tk-han-che}:
\begin{enumerate}
    \item Đưa khóa MoMo ra biến môi trường và thêm cơ chế chặn khởi động tương tự \mt{PayosProductionGuard} khi production dùng khóa sandbox.
    \item Thêm migration \mt{V4} bổ sung hai cột còn thiếu của bảng \mt{floors}, và một bước kiểm tra tự động đối chiếu lược đồ Flyway với các lớp thực thể.
    \item Dùng DTO riêng cho API tạo và sửa chính sách hủy; giới hạn danh sách người dùng theo chi nhánh của người gọi; bổ sung lọc theo loại không gian và thời gian hiệu lực của chính sách.
    \item Xóa quy tắc \mt{/h2-console/**} và bật lại ca kiểm thử \mt{h2ConsoleIsNotPublic}.
    \item Thiết lập CI (ví dụ GitHub Actions) chạy \mt{mvn test} và build giao diện cho mỗi thay đổi; thêm kiểm thử tích hợp dùng Testcontainers với PostgreSQL, trong đó có kiểm thử nhiều luồng đặt cùng một chỗ.
    \item Ghi nhật ký kiểm toán cho thao tác đổi vai trò, thu tiền mặt và check-in.
\end{enumerate}

\subsection{Dài hạn}
\begin{enumerate}
    \item \textbf{Bảo mật phiên đăng nhập:} chuyển token làm mới sang cookie \mt{HttpOnly}, thêm danh sách thu hồi token và giới hạn tần suất yêu cầu.
    \item \textbf{Chạy nhiều phiên bản máy chủ:} dùng ShedLock hoặc khóa trên PostgreSQL cho tác vụ định kỳ; chuyển bộ nhớ đệm dùng chung sang Redis.
    \item \textbf{Hoàn tiền tự động:} tích hợp API hoàn tiền của cổng thanh toán thay cho bước chuyển khoản thủ công.
    \item \textbf{Ứng dụng di động} cho khách hàng và cho nhân viên quét mã check-in.
    \item \textbf{Thiết bị tại chỗ:} kết nối khóa cửa hoặc cảm biến để tự động check-in và ghi nhận tình trạng sử dụng chỗ.
    \item \textbf{Cải thiện gợi ý đối tác:} thử nghiệm biểu diễn hồ sơ bằng vector nhúng (ví dụ lưu bằng pgvector) và so sánh với cách tính Jaccard hiện tại.
\end{enumerate}

\section{Kết luận}
\label{sec:tk-ket-luan}

Đồ án đã xây dựng được một hệ thống quản lý văn phòng chia sẻ nhiều chi nhánh chạy được trên môi trường đám mây, bao phủ các luồng chính: đặt chỗ, thanh toán trực tuyến, hủy và hoàn tiền, vận hành tại quầy, kết nối cộng đồng và báo cáo. Các quy tắc nghiệp vụ quan trọng, đặc biệt là chống đặt trùng lịch và xử lý webhook thanh toán, được kiểm thử bằng các ca tự động ở mức đơn vị và trên cơ sở dữ liệu.

Hệ thống chưa sẵn sàng cho vận hành thương mại: cần xử lý các điểm bảo mật tồn đọng (đặc biệt là cấu hình MoMo), bổ sung migration còn thiếu, kiểm thử tích hợp và CI trước khi đưa vào sử dụng thật. Qua đồ án, nhóm có thêm kinh nghiệm về thiết kế cơ sở dữ liệu có ràng buộc toàn vẹn, xử lý tương tranh, tích hợp cổng thanh toán và rà soát lỗi logic trong một hệ thống có nhiều vai trò người dùng.
